\section{Application to field data}
\label{eq:field_data}

We have applied the equivalent-layer technique to the Registro do Araguaia anomaly 
\citep{dutra_etal_2012, marangoni_mantovani_2013, dutra_etal_2014}, which is one the most prominent anomalies
over the NW portion of the Goiás Alkaline Province (GAP), Brazil \citep{junqueira-brod_etal_2005}. 
According to \citet{dutra_etal_2012}, there are no outcrops or any geological feature associated with 
this anomaly and the region is totally covered by quaternary sediments from the Araguaia River.
The GAP is a result of mafic-alkalic magmatism that occurred in the Late Cretaceous. Their intrusions 
penetrate the Bras{\'i}lia Belt basement and the Goi{\'a}s Magmatic Arch, as well as intrude and 
overlie basalts and sedimentary rocks from the northern border of the Paran{\'a} Basin.
\citet{dutra_etal_2014} consider that the source of Registro do Araguaia anomaly is one of the most 
important magma feeders for the GAP intrusions. 
They propose that this anomaly and its NE-SW lineament may be the main locus for magma conduit at the 
southern end of the Transbrasiliano Lineament.

Figure \ref{fig:Figure8} shows the location, terrain surface and Registro do Araguaia total-field anomaly.
The airborne data have N-S flight lines and E-W tight lines spaced from 0.5 km and 5.0 km, respectively. 
The airborne survey covered an area of approximately $58,834$ km$^{2}$ in the Goiás state, central Brazil,
from the beginning of July to the end of September 2004. 
In the study area, there are $480,830$ observation points with vertical coordinates varying from 
$-616.60$ m to $-358.41$ m, with average at $-426.43$ m and standard deviation $30.53$ m.

We have set a discrete equivalent layer on the horizontal surface $\mathcal{S}''$ at $z''_{0} = 291.10$ m, 
composed of punctual equivalent sources arranged on a regular grid of $N_{x} \times N_{y}$ points 
$\mathbf{r}''_{j} = (x''_{j}, y''_{j}, z''_{0})^{\top}$, $j \in \{1, \dots, M = N_{x} \, N_{y}\}$, with 
$N_{x} = 150$ and $N_{y} = 80$.

